\section{Формализация задачи как процесса обратного вывода конструктивной истории}

TODO2

В научно-техническом отношении задача проекта может быть формализована как задача обратного вывода конструктивной истории параметрической трехмерной модели. Пусть имеется последовательность состояний M0, M1, …, Mn, где каждый переход Mi→Mi+1 реализован некоторой операцией oi. Каждая операция обладает типом из словаря допустимых CAD-операций, параметрами, привязкой к геометрическим сущностям и местом в частично упорядоченной истории. Требуется по наблюдаемым данным Mk (где k >= i+1) восстановить отдельную операцию oi.

TODO3

Даже в таком компактном описании видно, что задача имеет гибридную природу. Пространство входов включает как дискретизированные поверхности STL, так и более точные топологические и параметрические описания STEP, а также текстовые или JSON-метаданные. Пространство выходов также неоднородно: тип операции является дискретной переменной; параметры включают непрерывные, целочисленные, булевы и категориальные величины; сегментация задает распределение меток по треугольникам, ребрам, граням или точкам; восстановление дерева построения требует последовательностной логики. Следовательно, простая редукция задачи к одной классификации или одной регрессии принципиально недостаточна.

Реализация проекта была построена в логике последовательности итерационных исследовательских циклов. Были рассмотрены два параллельных подхода: прямое предсказание операций и параметров по STL / текстового описания операций в формате Steps.txt, и альтернативный путь через классификацию и сегментирование каждой операции отдельно на основании STEP / текстового описания операций в формате JSON.

